% ── Referências ─────────────────────────────────────────────────────────────
\capitulopostext{REFERÊNCIAS}

\begin{singlespace}
\setlength{\parindent}{0pt}
\setlength{\parskip}{0pt}

ABNT -- Associação Brasileira de Normas Técnicas. \textbf{NBR 14724}: Informação e documentação. Trabalhos Acadêmicos - Apresentação. Rio de Janeiro: ABNT, 2002.

\vspace{\baselineskip}

BOYER, C. B.; UTA, C. M. \textbf{História da Matemática} [Trad. Helena Castro]. 3 ed. São Paulo: Blucher, 2012.

\vspace{\baselineskip}

D'AMBRÓSIO, U. \textbf{Educação Matemática:} da teoria à prática. 23. ed. Campinas: Papirus, 2012.

\vspace{\baselineskip}

KUBO, O.; BOTOMÉ, S. \textbf{Ensino e aprendizagem:} uma interação entre dois processos comportamentais. Interação, v.5, p.123-32, 2001.

\vspace{\baselineskip}

HART-DAVIS, A. \textbf{O Livro da Ciência.} 2. ed. São Paulo: Globo, 2016.

\vspace{\baselineskip}

PILETTI, C. \textbf{Didática geral.} São Paulo: Ática, 1995.

\vspace{\baselineskip}

RIBEIRO, J. L. P. Áreas e Proporções nas Superquadras de Brasília Usando o Google Maps.
\textbf{Revista do Professor de Matemática}. Rio de Janeiro, n. 92, p. 12-15, jan-abr. 2017.

\vspace{\baselineskip}

SEVERINO, A. J. \textbf{Metodologia do trabalho científico.} 22. ed. rev. e ampl. São Paulo: Cortez, 2002.

\end{singlespace}

\vspace{2\baselineskip}

\begin{singlespace}
\setlength{\parindent}{0pt}
O trabalho deverá ser redigido conforme recomendações das Diretrizes para confecção de teses e dissertações da Universidade de São Paulo (USP), disponíveis em:
\textless\url{http://www.teses.usp.br/index.php?option=com_content&view=article&id=52&Itemid=67}\textgreater.
Acesso em 24 jun.2021.
\end{singlespace}
